\documentclass[review,12pt]{elsarticle}
% Camera-ready (after acceptance): switch to
% \documentclass[final,3p,times,twocolumn]{elsarticle}

\usepackage{amsmath,amssymb,amsfonts}
\usepackage{algorithmic}
\usepackage{algorithm}
\usepackage{graphicx}
\usepackage{textcomp}
\usepackage{xcolor}
\usepackage{booktabs}
\usepackage{multirow}
\usepackage{tikz}
\usetikzlibrary{arrows.meta,fit,positioning,backgrounds}
\usepackage{hyperref}

\journal{Knowledge-Based Systems}

\begin{document}

\begin{frontmatter}

% ============================================================
% [BAN NHAP / DRAFT - Stage 2 cua quy trinh nghien cuu doc lap]
% Day la ban thao MOI, TACH RIENG khoi DH2A-KT (paper/main.tex).
% Chua co thuc nghiem that. Moi bang so lieu trong Section 7 la
% placeholder tuong minh, khong duoc dien so bia truoc khi chay
% harness that. Xem Section 9 "Reproducibility roadmap".
% ============================================================

\title{Agentic Write-Back in Knowledge Tracing: A Critic-Gated
Protocol for Closed-Loop Hypergraph Updates}

\author[utehy]{Tuan Dao Minh}
\author[utehy]{Khanh-Trinh Nguyen}
\author[utehy]{Duong Nguyen Tien}
\author[vnu]{Quoc Khanh Ngo}
\author[utehy]{Van-Hau Nguyen}
\author[vnu]{Le Hoang Son\corref{cor1}}
\ead{sonlh@vnu.edu.vn}
\cortext[cor1]{Corresponding author.}

\affiliation[utehy]{organization={Department of Quality Assurance and
Testing, Hung Yen University of Technology and Education},
city={Hung Yen},
country={Vietnam}}

\affiliation[vnu]{organization={Artificial Intelligence Research Center,
VNU Information Technology Institute, Vietnam National University},
city={Hanoi},
country={Vietnam}}

\begin{abstract}
Recent hypergraph-based knowledge tracing (KT) models capture high-order
exercise--concept correlations, and large language model (LLM) agents are
increasingly deployed as tutoring layers on top of such models. A
practice that has received little scrutiny is \emph{write-back}: an
agent, after diagnosing a learner and proposing a hint, records a new
structural edge (a session hyperedge) into the very graph the predictive
model will condition on in future sessions. This creates a closed loop
between an LLM's language-level reasoning and a numeric model's
structural input --- a loop that can either sharpen the graph with
session-grounded evidence or silently corrupt it with hallucinated
relations. We formalize this closed loop for hypergraph KT and ask three
questions: (i) does write-back measurably change downstream prediction
accuracy and calibration over successive sessions when the predictive
backbone is kept frozen; (ii) at what rate does the graph degrade when
the agent commits erroneous edges, and how much of that degradation a
mandatory verification gate (a ``Critic'') absorbs; and (iii) what
stopping rule should govern when further agent writes are no longer
worth their risk. We reuse an existing frozen hypergraph KT backbone, an
existing agent pipeline with a mandatory Critic gate, and an existing
write-back mechanism as infrastructure, and contribute three new
pieces: a sequential (round-by-round) evaluation harness that tracks
AUC and calibration drift as the graph accumulates agent-written edges;
a controlled corruption-injection protocol that measures the Critic's
block rate and precision/recall at catching synthetic bad edges,
together with a structural (DAG-compatibility) check adapted from a
prior leakage-audit protocol; and a drift-threshold stopping rule with a
bootstrap confidence interval. This paper reports the protocol, the
harness design, and the experimental plan; end-to-end numerical results
require a GPU host external to the drafting environment and are marked
as pending throughout (Section~7). We position the contribution
explicitly against two adjacent literatures --- LLM-agent-driven
knowledge graph enrichment, which has not evaluated impact on a
downstream predictive task, and dynamic/temporal hypergraph KT models,
which update graph structure algorithmically rather than through an
agent with an explicit accept/reject gate --- and against two companion
manuscripts by the same authors, from which this paper is deliberately
scoped apart.
\end{abstract}

\begin{keyword}
Knowledge tracing \sep hypergraph neural networks \sep large language
model agents \sep closed-loop learning \sep continual learning \sep
knowledge graph safety \sep human-AI collaborative systems
\end{keyword}

\end{frontmatter}

% ============================================================
\section{Introduction}
% ============================================================

Knowledge tracing (KT) predicts a learner's probability of answering a
future exercise correctly from their interaction history, and is the
core inference layer behind adaptive tutoring systems. Two trends have
converged in recent KT research. First, structural models that go beyond
pairwise concept graphs --- hypergraphs that connect an exercise to all
of its associated concepts at once --- have shown that high-order
exercise--concept correlations carry predictive signal that pairwise
graph neural network (GNN) models miss~\cite{lei2025hgkt,li2025dygas}.
Second, LLM agents are increasingly layered on top of KT backbones to
generate learner-facing explanations, hints, and next-step
recommendations, turning a purely numeric prediction into an
interpretable pedagogical interaction.

These two trends meet at a point that has not yet been studied
carefully: what happens when the agent is allowed to write back into the
graph the numeric backbone reads. Concretely, after a tutoring session
an agent may record a new hyperedge summarizing which concepts were
co-practiced, or which concept preceded which in the session --- a
plausible and, on the surface, beneficial way to keep the graph current
without waiting for a full retraining cycle. But the agent is a language
model operating under the same hallucination risk documented broadly in
the LLM literature: it can group concepts that were not actually related,
or infer a prerequisite relation that does not hold. Because the graph
is shared infrastructure that all future predictions condition on, a
single bad write does not stay local to one session; its effect
compounds every time the corrupted hyperedge participates in a future
forward pass.

Two concrete scenarios motivate the three research questions of this
paper.

\emph{Scenario 1 (does write-back help).} A learner struggles with two
exercises that share a concept the graph did not previously link at the
hyperedge level (e.g. a course-specific skill not present in the
original curriculum graph). After the session, the agent writes a
hyperedge connecting the exercises' concepts. On the learner's next
session, does the backbone's frozen weights, now reading a richer graph,
produce a measurably better probability estimate and better-calibrated
confidence than they would have on the old graph -- or is the effect too
small to detect above noise, in which case the added structural
complexity is not paying for itself?

\emph{Scenario 2 (does write-back hurt).} The agent, prompted by a
surface-level lexical similarity between two exercise stems, writes a
hyperedge linking concepts that are not actually related. The Critic
gate --- an existing mandatory verification step that flags contradiction
with the backbone's own probability estimate --- may or may not catch
this, since a single spurious edge does not necessarily produce an
immediately visible contradiction. If it is not caught, the corrupted
edge persists and participates in every subsequent hypergraph
convolution. How many further sessions elapse, and how many further
writes accumulate, before the corruption becomes visible in downstream
accuracy -- and is that too late to be useful as a safety signal?

Both scenarios point to the same underlying object of study: the
\emph{trajectory} of a hypergraph KT system's downstream performance as
a function of the number of agent-committed writes, under a verification
gate whose block rate and error rate are themselves measurable
quantities. Existing work does not provide this trajectory. Multi-agent
frameworks for LLM-driven knowledge graph enrichment, such as
KARMA~\cite{lu2025karma}, evaluate enrichment quality by LLM-verified
correctness on the enrichment task itself (entity and relation
extraction from text), not by the effect of the enriched graph on a
downstream predictive model's accuracy over time. Dynamic and temporal
hypergraph KT models, such as the Temporal Hypergraph Memory
Network~\cite{mohammadi2025thmn}, DyGAS~\cite{li2025dygas}, and
HGKT~\cite{lei2025hgkt}, update graph structure or gating weights through
a fixed algorithm learned once at training time; none of them involves an
LLM agent making a discrete accept/reject decision about a candidate
structural edit after deployment, and none reports a degradation
trajectory under adversarial or noisy edits. General continual-learning
work on concept drift~\cite{ashrafee2025amr} addresses distributional
drift in the input stream, not drift deliberately introduced by the
system's own agent into its own structural memory. Multi-agent
self-checking frameworks for LLM hallucination, such as
MARCH~\cite{li2026march}, provide a verification pattern (a
Solver/Proposer/Checker pipeline with an explicit accept/reject verdict)
that is the closest structural analogue to the Critic gate studied here,
but have been evaluated on retrieval-augmented question answering, not
on structural edits to a numeric model's input graph.

This paper's contribution is therefore not a new predictive backbone and
not a new agent architecture; both already exist as infrastructure
(Section~4). The contribution is a \emph{protocol and harness} for
answering the three research questions above, and a first experimental
plan for running it. We are explicit that, at the time of this draft, the
harness has been designed and specified but not yet executed at scale;
Section~7 reports the exact status of every quantity the harness is
meant to produce, using the same placeholder convention as our companion
manuscripts (Section~2) rather than filling tables with estimated or
illustrative numbers.

\subsection{Relationship to companion manuscripts by the same authors}
\label{sec:relationship}

This manuscript reuses code and infrastructure from two companion works
by the same author group and is explicitly scoped apart from both, to
avoid redundant contribution claims.

\cite{daominh2026p0} (submitted to Applied Intelligence) is a
leakage-controlled concept-graph construction and cold-start diagnostic
\emph{protocol}: learner-based temporal splits, train-only prerequisite
inference, DAG validation with cycle pruning, and a DAG Disruption Rate
(DDR) diagnostic, applied once, statically, before any model is trained.
It does not involve an agent, does not run at inference time, and does
not study anything that happens after deployment.

Our companion manuscript on DH$^2$A-KT (submitted to this journal,
concurrently) describes a full two-tier system: a frozen hypergraph KT
backbone (Tier~1) and an agentic explanation/tutoring pilot (Tier~2)
that includes the write-back mechanism as one architectural component
among several (alongside a causal hint-effect estimator and an
explanation-faithfulness evaluation). DH$^2$A-KT reports a single
end-to-end system and treats write-back descriptively, as part of the
Tier~2 loop; it does not measure a downstream-performance trajectory
across multiple rounds of writes, does not run a controlled
corruption-injection study, and does not propose a stopping rule.

This paper isolates write-back as the sole object of study, asks
questions neither companion paper answers, and can be read, reviewed,
and (if applicable) accepted independently of either. Table~\ref{tab:scope-comparison}
summarizes the scope boundaries.

\begin{table}[!t]
\caption{Scope boundaries versus the two companion manuscripts. A
checkmark means the item is a central, evaluated contribution of that
paper; a dash means the item is absent or only mentioned in passing.}
\label{tab:scope-comparison}
\centering
\small
\begin{tabular}{lccc}
\toprule
& P0 (APIN, \cite{daominh2026p0}) & DH$^2$A-KT (companion, this journal) & This paper \\
\midrule
Static leakage/DAG audit protocol & \checkmark & reused (frozen) & reused (frozen) \\
Frozen hypergraph KT backbone (Tier~1) & --- & \checkmark & reused (frozen) \\
Agentic explanation pilot (Tier~2) & --- & \checkmark & reused (infrastructure) \\
Causal hint-effect (ATE/IPW) estimation & --- & \checkmark & --- \\
Explanation faithfulness (KC-Jaccard, IG control) & --- & \checkmark & --- \\
Multi-round write-back trajectory (AUC/calibration drift) & --- & --- & \checkmark \\
Controlled corruption injection into write-back & --- & --- & \checkmark \\
Critic block-rate / precision-recall analysis & --- & mentioned & \checkmark \\
Drift-based stopping rule & --- & --- & \checkmark \\
\bottomrule
\end{tabular}
\end{table}

\subsection{Contributions}

\begin{enumerate}
\item A closed-loop evaluation harness that runs the existing Tier~2
agent pipeline over a sequence of sessions, holding the Tier~1 backbone
frozen, and records AUC and expected calibration error (ECE) after every
committed write (Section~5.1).
\item A controlled corruption-injection protocol that inserts
synthetic, known-bad candidate hyperedges at a specified rate, extends
the existing Critic gate with a structural DAG-compatibility check
reused from~\cite{daominh2026p0}'s cycle-pruning logic, and reports the
gate's block rate, precision, and recall against the injected ground
truth (Section~5.2).
\item A drift-threshold stopping rule with a bootstrap confidence
interval that flags when cumulative structural drift should pause
further agent writes (Section~5.3).
\item An explicit, auditable status report of every quantity above at
the current stage of implementation (Section~7), following the
disclosure convention already used in our companion manuscript's
faithfulness tables.
\end{enumerate}

% ============================================================
\section{Related Work}
% ============================================================

\subsection{Hypergraph and dynamic-graph knowledge tracing}

Hypergraph KT models connect an exercise to the full set of concepts it
touches in a single hyperedge, rather than decomposing that relation into
pairwise edges. HGKT~\cite{lei2025hgkt} uses a two-layer hypergraph
convolution together with a learning-layer/forgetting-layer dual-gating
mechanism to balance how much new evidence and how much decay affect the
knowledge state; the dual gate is internal to a single forward pass over
a graph that is otherwise static after training. DyGAS~\cite{li2025dygas}
augments sequence modeling with a dynamically evolving graph, and the
Temporal Hypergraph Memory Network~\cite{mohammadi2025thmn} models
concept interactions through bidirectional message passing over a
temporal hypergraph with an adaptive scaling mechanism. In all three,
graph evolution (if present) is computed by a fixed, differentiable
update rule learned end-to-end; there is no point at which an external
agent proposes a discrete structural edit that a separate mechanism can
accept or reject. \textbf{Note for the authors:} HGKT's dual-gating
mechanism is close enough in spirit to the \texttt{DualGatedUpdate}
module already used in our Tier~1 backbone (forget/input gates over the
concatenation of previous state and new evidence) that the eventual
camera-ready version must state explicitly, with an equation-level
comparison, how the two differ (HGKT gates a learning/forgetting
\emph{decomposition} of the update; our backbone gates a
forget/candidate decomposition directly) --- this is a citation
placed correctly, but the differentiation needs to be sharper before
submission.

\subsection{LLM agents and dynamic knowledge graphs}

KARMA~\cite{lu2025karma} is the closest existing precedent for an LLM
agent that writes into a knowledge graph under a verification step: nine
collaborative agents, including a dedicated Conflict Resolution Agent
that runs an LLM-based debate mechanism, enrich a biomedical knowledge
graph from unstructured text, reporting 83.1\% LLM-verified correctness
on newly added triples and an 18.6\% reduction in conflicting edges.
KARMA's evaluation, however, is intrinsic to the enrichment task itself
(is the new triple correct according to the LLM's own judgment); it does
not measure whether the enriched graph changes the output of a separate
downstream predictive model, and it is not applied in a setting where the
graph is read by a frozen, previously trained scoring function on a
recurring basis. Agent-memory frameworks built on temporal or evolving
knowledge graphs (e.g. graph-structured agent memory substrates) treat
the graph as a store of facts about the world or the conversation, not as
the direct numeric input to a supervised model whose accuracy can be
measured against ground-truth labels; this is a structural difference
that lets the present paper use a much sharper success criterion (AUC,
ECE) than ``was this fact judged correct.''

\subsection{Verification gates against agent hallucination}

MARCH~\cite{li2026march} formalizes a Solver/Proposer/Checker pipeline in
which the Checker validates claim-level propositions against retrieved
evidence in isolation from the Solver's own output, producing an
explicit verdict. The general pattern --- decompose a candidate output
into checkable units, verify each independently of the agent that
produced it, and emit an explicit accept/reject decision --- is the same
pattern the existing Critic gate in our Tier~2 pipeline already
implements for explanation contradiction. This paper's second
contribution (Section~5.2) extends that same Critic to a second,
structural checking task (DAG-compatibility of a candidate hyperedge)
that MARCH's text-claim setting has no analogue for, and reports metrics
(block rate, precision, recall against known-injected corruption) that
verifier papers in the text-generation literature do not typically
report because they lack a ground-truth corruption channel; our
controlled-injection design (Section~5.2) is what supplies that ground
truth.

\subsection{Continual learning under concept drift}

Recent continual-learning work explicitly separates virtual drift
(handled by classical continual learning, which mitigates forgetting)
from real, representation-level drift, proposing detection modules
coupled with adaptive memory realignment~\cite{ashrafee2025amr}. This
literature treats drift as something that happens to the input
distribution and must be detected and compensated for; it does not
consider a setting where the system's own agent is the source of the
drift by design, nor a setting where a verification gate can intercept
the drift-causing event (a candidate graph edit) before it is committed,
rather than only after its effect appears in the data stream. Framing
agent write-back as a controllable, gate-able source of drift, rather
than an environmental nuisance to be detected after the fact, is the gap
this paper's methodology (Section~5) is built to close.

% ============================================================
\section{Problem Formulation}
% ============================================================

Let $G_0 = (V, E_0)$ denote the hyperedge-augmented concept graph at
deployment time, where $V$ is the set of concepts and $E_0$ the audited
hyperedges inherited, unchanged, from~\cite{daominh2026p0}. Let
$f_\theta$ denote the Tier~1 hypergraph KT backbone, trained once on
$G_0$ and \emph{frozen} thereafter (weights $\theta$ do not update after
deployment; this mirrors the frozen-backbone design already used for the
Tier~1/Tier~2 boundary in DH$^2$A-KT). A \emph{round} $t = 1, 2, \dots$
corresponds to one completed tutoring session in which the existing
Tier~2 pipeline runs to completion: the Diagnostician produces a
grounded explanation from $f_\theta$'s frozen output and the learner's
history; the Tutor/Hint agent proposes a next hint and, from the
session's interaction pattern, drafts a candidate hyperedge
$e_t^{\text{cand}}$ (for example, linking the concepts practiced together
in that session); the Critic evaluates $e_t^{\text{cand}}$ and emits
$\text{APPROVE}$ or $\text{REJECT}$.

The graph evolves as
\[
G_t =
\begin{cases}
G_{t-1} \cup \{e_t^{\text{cand}}\} & \text{if Critic APPROVEs} \\
G_{t-1} & \text{if Critic REJECTs.}
\end{cases}
\]

Because $\theta$ is frozen, $G_t$ affects future predictions only through
$f_\theta$'s hypergraph convolution operating over a different input
structure, not through any weight update -- this is a deliberate
simplification that isolates the effect of \emph{structural} change from
the effect of \emph{parametric} change, and is consistent with the
frozen-Tier1/agentic-Tier2 boundary already established in the companion
manuscript. Let $\mathcal{D}_t^{\text{eval}}$ be a held-out set of future
interactions not seen by round $t$. We define:

\[
\Delta\text{AUC}_t = \text{AUC}\big(f_\theta(\cdot \mid G_t),\ \mathcal{D}_t^{\text{eval}}\big) - \text{AUC}\big(f_\theta(\cdot \mid G_0),\ \mathcal{D}_t^{\text{eval}}\big)
\]

\[
\Delta\text{ECE}_t = \text{ECE}\big(f_\theta(\cdot \mid G_t),\ \mathcal{D}_t^{\text{eval}}\big) - \text{ECE}\big(f_\theta(\cdot \mid G_0),\ \mathcal{D}_t^{\text{eval}}\big)
\]

Research question (i) asks whether $\Delta\text{AUC}_t$ and
$\Delta\text{ECE}_t$ move in the beneficial direction, and by how much,
under normal (non-adversarial) agent behavior. Research question (ii)
asks how $\Delta\text{AUC}_t$ and $\Delta\text{ECE}_t$ evolve when a
known fraction $p_{\text{inject}}$ of candidate hyperedges are
synthetically corrupted before reaching the Critic, and what fraction of
those the Critic blocks (Section~5.2). Research question (iii) asks for
a decision rule $\tau(\{\Delta\text{AUC}_1, \dots, \Delta\text{AUC}_t\})
\to \{\text{continue}, \text{pause}\}$ that a deployed system could use
without access to future ground truth (Section~5.3).

% ============================================================
\section{Reused Infrastructure}
% ============================================================

Three components are reused, unchanged, from existing code in this
project and are treated as fixed infrastructure rather than as
contributions of this paper.

\textbf{Tier~1 backbone.} The frozen hypergraph KT model: a concept
embedding table, a per-hyperedge-kind stack of hypergraph convolution
layers, a dual-gated recurrent update over the fused concept/response
representation, and a linear output head producing $P(\text{correct})$.
Training procedure, hyperparameters, and the auxiliary graph-sensitivity
loss are exactly as reported in the companion DH$^2$A-KT manuscript and
are not modified here.

\textbf{Tier~2 pipeline and Critic gate.} The existing Diagnostician /
Critic / Tutor-Hint agent loop (\texttt{dh2a\_kt/agents/diagnostician.py},
\texttt{dh2a\_kt/tier2/pilot.py}), including the mandatory Critic
verification step that already flags contradiction between an agent's
explanation and the Tier~1 probability estimate, and the existing
\texttt{ablation}/\texttt{skip\_critic} parameters that already allow
disabling the Critic for controlled comparison
(\texttt{scripts/04\_run\_tier2\_pilot.py --skip-critic}). This paper
does not modify the Critic's existing contradiction check; it adds a
second, structural check (Section~5.2) alongside it.

\textbf{Write-back mechanism.} The existing session-hyperedge write-back
step, by which an approved candidate hyperedge is appended to the graph
consumed by subsequent Tier~1 forward passes.

% ============================================================
\section{Method: Closed-Loop Evaluation Protocol}
% ============================================================

\subsection{Sequential drift-tracking harness (RQ i)}

The harness runs $T$ sessions in sequence (rather than the single
end-to-end pilot reported in the companion manuscript) under four
conditions, held constant across a session order fixed by a random seed:

\begin{enumerate}
\item \textbf{Static} --- Tier~2 disabled; $G_t = G_0$ for all $t$. Control
condition isolating any drift not caused by write-back at all (e.g.
evaluation-set composition effects).
\item \textbf{Write-back, no Critic} --- \texttt{skip\_critic=True}; every
candidate hyperedge is committed regardless of quality. Upper bound on
risk.
\item \textbf{Write-back, Critic (contradiction check only)} --- the
existing default behavior.
\item \textbf{Write-back, Critic (contradiction + DAG-compatibility)} ---
the extended gate proposed in Section~5.2.
\end{enumerate}

After every round $t$, $\Delta\text{AUC}_t$ and $\Delta\text{ECE}_t$ are
computed on a held-out evaluation slice and logged; the four conditions
are compared as trajectories over $t = 1, \dots, T$, not as single
end-point numbers, since the research question is about the shape of the
trajectory (monotonic improvement, plateau, or eventual degradation) as
much as its endpoint.

\subsection{Controlled corruption injection and Critic evaluation (RQ ii)}

At a specified rate $p_{\text{inject}}$, the harness substitutes a
synthetically corrupted candidate hyperedge for the agent's real one
before it reaches the Critic. Two corruption types are used, both
constructible without new data collection: (a) \emph{concept
swap} --- replace one endpoint concept with a randomly sampled,
unrelated concept, simulating a lexical-similarity hallucination; (b)
\emph{direction/cycle violation} --- construct a candidate edge that
would create a cycle in the prerequisite sub-graph, reusing the same
cycle-detection routine already validated for the static audit
in~\cite{daominh2026p0}. Because the ground-truth corruption label is
known by construction, the Critic's decision on each injected item can be
scored as a binary classifier: block rate (recall on corrupted items),
false-block rate (1 -- specificity on genuine items), and precision.
This is the key methodological difference from KARMA's evaluation
(Section~2.2): KARMA reports what fraction of the agent's \emph{own}
output an LLM judge considered correct, whereas this protocol reports
what fraction of \emph{known-corrupted} synthetic input a separate gate
catches, which is a stricter and independently verifiable signal.

\subsection{Drift-threshold stopping rule (RQ iii)}

We define a rolling-window drift statistic
$s_t = \Delta\text{AUC}_t - \Delta\text{AUC}_{t-w}$ over a window of $w$
rounds, and a bootstrap confidence interval for $s_t$ constructed by
resampling evaluation examples within the window (reusing the bootstrap
machinery already implemented for AUC confidence intervals
in~\cite{daominh2026p0}). The rule proposes pausing further writes when
$s_t$'s bootstrap CI excludes zero in the negative direction for
$k$ consecutive rounds, where $w$ and $k$ are harness parameters swept
in the experimental design (Section~6) rather than fixed a priori.

% ============================================================
\section{Experimental Design}
% ============================================================

\textbf{Data.} FoundationalASSIST is the primary corpus, since it is the
same dataset already used for the companion manuscript's Tier~2 pilot and
therefore requires no new access approval; XES3G5M is used as a secondary
corpus for the Static and no-Critic conditions only (conditions that do
not require live LLM sessions and can therefore run at larger $T$).

\textbf{Backbone and LLM.} Frozen Tier~1 weights from the companion
manuscript's checkpoint (no retraining). Tier~2 agent runs on Qwen2.5-7B
via Ollama, matching the companion manuscript's configuration, so that
any difference in trajectory is attributable to the closed-loop protocol
rather than to a different underlying LLM.

\textbf{Session budget.} $T = 200$ sequential sessions per condition per
seed, 3 seeds, matching the order of magnitude of the companion
manuscript's 500-sample single-round pilot while keeping total LLM calls
tractable on a single RTX~3090 host (the same hardware budget documented
in \texttt{docs/kbs-gpu-runbook.md} for the companion manuscript).

\textbf{Injection rate sweep.} $p_{\text{inject}} \in \{0, 0.05, 0.10,
0.20\}$ for the corruption-injection arm (Section~5.2), matching the
order of magnitude of injection rates already used for the Tier~1
manipulation check in the companion manuscript ($p=0.9$ node-drop is a
different, more aggressive perturbation intended to stress-test the
backbone directly; the rates here are chosen to be plausible per-session
hallucination rates rather than a worst-case stress test).

\textbf{Metrics reported per condition, per round.} $\Delta\text{AUC}_t$,
$\Delta\text{ECE}_t$, cumulative Critic block rate, block-rate precision
and recall against injected corruption, and the stopping rule's trigger
round (if any) under the $w,k$ sweep.

% ============================================================
\section{Results (status: pending -- harness not yet executed)}
% ============================================================

\textbf{This section is deliberately left as a set of table shells.}
Unlike the companion manuscript's faithfulness harness, which already has
a stub-LLM smoke run validating pipeline wiring at the time of this
draft, the sequential drift-tracking harness described in Section~5 has
been specified but not yet implemented in code. No numbers, stub or
otherwise, exist for it. Filling the tables below with placeholder,
illustrative, or literature-typical numbers before the harness is run
would misrepresent the current state of the work; every cell is therefore
marked \texttt{---} and must not be edited except by inserting a value
produced by an actual harness run, with the run's script, seed, and log
path recorded in the caption.

\begin{table}[!t]
\caption{RQ (i): $\Delta$AUC and $\Delta$ECE trajectory summary
(endpoint at $t{=}200$, mean $\pm$ SD over 3 seeds). \textbf{Status:
harness not yet implemented; no run has occurred.} Populate only after
\texttt{scripts/2X\_run\_closedloop\_harness.py} (to be written) executes
on FoundationalASSIST per Section~6.}
\label{tab:rq1-placeholder}
\centering
\small
\begin{tabular}{lcc}
\toprule
Condition & $\Delta\text{AUC}_{200}$ & $\Delta\text{ECE}_{200}$ \\
\midrule
Static (no write-back) & --- & --- \\
Write-back, no Critic & --- & --- \\
Write-back, Critic (contradiction only) & --- & --- \\
Write-back, Critic (contradiction + DAG) & --- & --- \\
\bottomrule
\end{tabular}
\end{table}

\begin{table}[!t]
\caption{RQ (ii): Critic performance against synthetic corruption
injection, by injection rate. \textbf{Status: pending, same harness.}}
\label{tab:rq2-placeholder}
\centering
\small
\begin{tabular}{lcccc}
\toprule
$p_{\text{inject}}$ & Block rate (recall) & Precision & False-block rate & Residual $\Delta$AUC \\
\midrule
0.05 & --- & --- & --- & --- \\
0.10 & --- & --- & --- & --- \\
0.20 & --- & --- & --- & --- \\
\bottomrule
\end{tabular}
\end{table}

\begin{table}[!t]
\caption{RQ (iii): stopping-rule trigger round under the $(w,k)$ sweep.
\textbf{Status: pending, same harness.}}
\label{tab:rq3-placeholder}
\centering
\small
\begin{tabular}{ccc}
\toprule
Window $w$ & Consecutive rounds $k$ & Trigger round (mean over seeds) \\
\midrule
10 & 2 & --- \\
10 & 3 & --- \\
20 & 2 & --- \\
20 & 3 & --- \\
\bottomrule
\end{tabular}
\end{table}

% ============================================================
\section{Discussion}
% ============================================================

Even without numerical results, the protocol itself yields two claims
that can be stated now. First, the four-condition design
(Section~5.1) is directly comparable across conditions because three of
its four arms already exist as executable code paths in the reused
infrastructure (Static and no-Critic are direct consequences of existing
CLI flags; only the contradiction+DAG condition requires new code), which
substantially lowers the implementation risk relative to a from-scratch
continual-learning study. Second, the corruption-injection design
(Section~5.2) supplies a ground-truth label that neither KARMA's
self-judged correctness nor the companion manuscript's single-round
Critic pass rate provide, and is what would let a revised version of this
paper report precision/recall for the Critic rather than only a pass/fail
count -- a meaningfully stronger evaluation than exists anywhere in the
closed-loop-agent-and-KT intersection at the time of writing, to the best
of our knowledge.

Once run, three outcomes are plausible and each has a different
publishable story. If $\Delta\text{AUC}_t$ under the Critic-gated
conditions is positive and significantly above the no-Critic condition,
the paper's claim is that gated write-back is a safe, low-cost way to
keep a frozen KT backbone's effective input current between retraining
cycles. If $\Delta\text{AUC}_t$ is statistically indistinguishable from
zero across all conditions, the paper's claim shifts to a negative but
still useful result: that session-level write-back, at the rate and
granularity tested, is not a reliable route to better prediction, and
that the Critic's value lies entirely in bounding downside risk
(Table~\ref{tab:rq2-placeholder}) rather than in enabling upside. If
$\Delta\text{AUC}_t$ degrades even under the full gate, the paper's claim
is that the current single-step contradiction check is an insufficient
safeguard for structural edits, motivating the DAG-compatibility
extension as necessary rather than optional -- in which case
Table~\ref{tab:rq2-placeholder}'s precision/recall breakdown would be the
paper's central table. All three outcomes are reportable; none require
the protocol to be redesigned after the fact, which is the property a
pre-registered-style harness is meant to provide.

% ============================================================
\section{Limitations}
% ============================================================

The frozen-backbone simplification (Section~3) means this study does not
address the case where periodic retraining is combined with write-back;
a system that retrains on the accumulated agent-written graph could show
a different, possibly larger, drift trajectory than the one this harness
measures. The corruption types in Section~5.2 (concept swap, cycle
violation) are two plausible failure modes among many an LLM agent could
produce; they are chosen because they are constructible without new
annotation and reuse existing cycle-detection code, not because they are
known to be the most common real failure mode -- that remains an open
empirical question the harness itself could eventually inform by
comparing injected-corruption trajectories against organically-observed
agent errors. $T=200$ sessions is a budget constraint tied to a single
consumer GPU host, not a principled stopping point for the trajectory
itself; longer horizons may reveal effects (slow compounding drift, in
particular) that this budget cannot rule out.

\section*{CRediT authorship contribution statement}
% [CAN TAC GIA XAC NHAN LAI] -- vai tro duoi day la suy luan hop ly dua
% tren thu tu tac gia, giong cach lam da dung cho ban thao DH2A-KT KBS,
% KHONG phai do tung dong tac gia tu khai bao. Phai ra soat truoc khi nop.
\textbf{Tuan Dao Minh:} Methodology, Software, Validation, Formal
analysis, Writing -- original draft.
\textbf{Khanh-Trinh Nguyen:} Validation, Writing -- review \& editing.
\textbf{Duong Nguyen Tien:} Validation, Writing -- review \& editing.
\textbf{Quoc Khanh Ngo:} Investigation, Writing -- review \& editing.
\textbf{Van-Hau Nguyen:} Resources.
\textbf{Le Hoang Son:} Conceptualization, Supervision, Project
administration, Writing -- review \& editing.

\section*{Declaration of competing interest}
% [CAN TAC GIA XAC NHAN]
The authors declare that they have no known competing financial
interests or personal relationships that could have appeared to
influence the work reported in this paper.

\section*{Data availability}
% [CAN TAC GIA XAC NHAN]
The closed-loop evaluation harness described in Section~5 will be
released alongside the companion DH$^2$A-KT repository once implemented.
FoundationalASSIST~\cite{foundationalassist2026} must be obtained
directly from Hugging Face under its Responsible Use Guidelines.

\section*{Acknowledgment}
This study is designed to run on the same single consumer-grade
NVIDIA RTX~3090 (24GB) budget already validated for the companion
manuscript's Tier~2 pilot.

\bibliographystyle{elsarticle-num}
\bibliography{references}

\end{document}
